% Fixture: proofs longer than three pages, citations wrapped to the start of a line, and end marks that sit far from
% the text: after a short last line of a wide page, and alone on their row under a display. Compile with: tectonic long-proofs.tex
\documentclass[11pt]{article}
\usepackage{amsmath,amssymb,amsthm}
\usepackage[margin=0.9in]{geometry}
\newtheorem{theorem}{Theorem}
\newtheorem{lemma}[theorem]{Lemma}
\newtheorem{proposition}[theorem]{Proposition}
\theoremstyle{remark}
\newtheorem{remark}{Remark}
\newcommand{\filler}{The argument continues by expanding each term and collecting the powers of the step size, which
leaves a sum whose coefficients are bounded uniformly in the parameters. Every estimate below uses only the bound
stated at the start of the proof, so the constants do not depend on the stiff part of the problem. Repeating this
for the remaining terms gives the same conclusion, and the remainders are absorbed into the final constant.\par}
\newcommand{\fillers}{\filler\filler\filler\filler\filler\filler\filler\filler}
\begin{document}

\section{Expansions}
We work with a family of bounded operators and with a step size that is smaller than one throughout.

\begin{lemma}\label{lem:bound}
The resolvent is bounded by a constant that does not depend on the step size.
\end{lemma}
\begin{proof}
This is the usual estimate, and it needs no assumption beyond the ones already stated.
\end{proof}

\begin{theorem}\label{thm:expansion}
The error has an expansion in powers of the step size.
\end{theorem}
\begin{proof}
\fillers\fillers
We substitute the first expansion into the second one and bound the resolvent. The terms collect as before, and we use\linebreak
Lemma~\ref{lem:bound}:
\[ \|x\| \le C h. \]
\fillers\fillers
The remainder is of the size that the first estimate in this proof gives for $\|x\|, \|y\|, \|z\| \le C h$ by\linebreak
Lemma~\ref{lem:bound}(a), and the constant stays the same as before.
\fillers\fillers
\fillers\fillers
Collecting everything gives the expansion. This completes the induction on the order of the
expansion.\linebreak
All done.
\end{proof}

The expansion has a first consequence, which the next statement records for later use in the following
\begin{remark}\label{rem:first}
The constants above are uniform in the step size.
\end{remark}

\begin{proposition}\label{prop:display}
The two errors agree to first order.
\end{proposition}
\begin{proof}
Subtract the two expansions, which by the argument in Theorem~\ref{thm:expansion} leaves the difference
\[ e_1 - e_2 = \mathcal{O}(h^2). \]
\end{proof}

The sentence after this proof does not belong to it, and it mentions Lemma~\ref{lem:bound}.

\end{document}
